%!TEX root = ./main.tex

\section[Intro]{Introduction}

% SECTION TIMING: 0:00--2:20 (1 frame).

% Teaching note (140 s): open with the picture and the question; do not answer it
% yet. Take one or two guesses. Receiver A needs the video itself; receiver B only
% needs to know where the pedestrians are. The whole lecture lives in the gap between
% those two receivers. Then reveal the route one line at a time.
\begin{frame}{One video stream, two receivers}
  \vspace{0.6em}
  \centering
  \begin{tikzpicture}[scale=0.85,transform shape]
    \node[netnode,minimum width=2.6cm,minimum height=1.1cm] (cam) at (0,0)
      {Street camera\\[-2pt]{\scriptsize\mdseries live video}};
    \node[netnode,minimum width=4.4cm,minimum height=1.0cm] (a) at (6.4,0.9)
      {A: a person\\[-2pt]{\scriptsize\mdseries wants to \textit{watch} the video}};
    \node[corenode,minimum width=4.4cm,minimum height=1.0cm] (b) at (6.4,-0.9)
      {B: an AI at the edge\\[-2pt]{\scriptsize\mdseries asks ``any pedestrians?''}};
    \draw[flow] (cam.east) -- (a.west);
    \draw[flow] (cam.east) -- (b.west);
    \node[lbl] at (2.65,0) {wireless};
  \end{tikzpicture}

  \vspace{0.5em}
  \askat{2-}{Do A and B need the same data?}

  \vspace{0.6em}
  \raggedright
  \begin{enumerate}[<+(2)->]\itemsep4pt
    \item Today: copy every bit of every frame
    \item Semantic communication: send the \textbf{meaning}
    \item Version 1: \textbf{rebuild} the video
    \item Version 2: just get the \textbf{task} done --- no video rebuilt
    \item Version 3: keep today's pipe, spend bits by \textbf{meaning}
  \end{enumerate}
\end{frame}
